% Source-aligned English for printed p.177.
\SourcePage{177}

\LayerTarget{tohoku:corollary:3.8.1:4:p2}
\noindent\textit{$(U_1,\ldots\mathpunct{.}U_k)$ of open sets in $\mathbf U$, their
intersection $U$ satisfies $\check H^i(U,F)=\{0\}$ for $i>0$. Then also
$H^i(U,F)=\{0\}$, and for every open subset $V$ of $X$, the natural
homomorphism $\check H^i(V,F)\to H^i(V,F)$ is an isomorphism.}

\LayerTarget{tohoku:proof:3.8.1:corollary-4:fixed-cover}
It is enough to prove $H^i(U,F)=\{0\}$, since $V$ has arbitrarily fine
coverings $\mathbf R$ by open sets in $\mathbf U$, and Corollary 1 shows that
for such an $\mathbf R$ the homomorphism
$H^i(R,F)\to H^i(V,F)$ is an isomorphism. At the same time,
since $\mathbf R$ is arbitrarily fine, this proves that
$\check H^i(V,F)\to H^i(V,F)$ is an isomorphism.
\LayerTarget{tohoku:proof:3.8.1:corollary-4:induction}
--- To prove $H^i(U,F)=\{0\}$, we prove by induction on $n$ that
$H^i(U,F)=\{0\}$ for $0<i\leqq n$ and every $U$ of the indicated form. This is
trivial for $n=0$. Suppose $n\geqq 1$ and that the assertion has been proved
for $n'=n-1$. There are arbitrarily fine coverings $\mathbf R$ of $U$ by open
sets in $\mathbf U$; for such an $\mathbf R$ we have
\LayerTarget{tohoku:eq:corollary-4-cover-complex-vanishing}
$C(\mathbf R,H^q(F))=0$ for $0<q<n$
by the induction hypothesis, and hence, a fortiori,
\LayerTarget{tohoku:eq:corollary-4-cover-cohomology-vanishing}
$H^p(R,H^q(F))=0$ for such $q$.
It follows that
\LayerTarget{tohoku:eq:corollary-4-cech-limit-vanishing}
$\check H^p(U,H^q(F))=0$ for such $q$.
By the observation preceding Corollary 4, this implies
\LayerTarget{tohoku:eq:corollary-4-induction-conclusion}
$H^n(U,F)=\check H^n(U,F)$,
which is zero.

\LayerTarget{tohoku:application:affine-zariski-coherent-sheaf}
Corollary 4 applies, for example, when $X$ is an algebraic variety with its
Zariski topology, $\mathbf U$ is the collection of affine open subsets of
$X$, and $F$ is a coherent algebraic sheaf on $X$ [15]. Indeed, by [15], the
affine open subsets form a basis for the topology of $X$, and the intersection
of two affine open subsets is affine. If $U$ is affine open, then
$\check H^i(U,F)=0$ for $i>0$. Thus
\LayerTarget{tohoku:eq:affine-cech-derived-comparison}
$H^i(X,F)=\check H^i(X,F)$.
Moreover, Corollary 1 above shows that the $H^i(X,F)$ can be computed using an
arbitrarily chosen covering of $X$ by affine open subsets.

\LayerTarget{tohoku:remark:closed-cover-leray-spectral-sequences}
\textsc{Remark.} There are cases other than that of Theorem 3.8.1 in
which the Leray spectral sequence is valid. The best-known one is a
\textit{locally finite covering of a paracompact $X$ by closed subsets} (the
case considered by Leray); it is most simply treated as above for open coverings, since the restriction of a soft sheaf to a closed subset is still
soft (replacing Proposition 3.1.3). Another case, obtained by Godement
by a different method, is a \textit{finite covering of $X$ by closed subsets}
(without a paracompactness hypothesis). When both cases apply simultaneously,
the two resulting spectral sequences coincide, fortunately.

\LayerTarget{tohoku:example:cech-H2-comparison-failure:p1}
\textsc{An example.} To conclude this section, we shall give a simple
example in which the monomorphism
$\check H^2(X,F)\to H^2(X,F)$ is not an isomorphism, and in fact
$\check H^2(X,F)=0$ while $H^2(X,F)\neq 0$. Since Corollary 2 to Theorem 3.8.1
gives an exact sequence
\LayerTarget{tohoku:eq:cech-H2-comparison-exact-sequence}
\[
0\to\check H^2(X,F)\to H^2(X,F)
\to\check H^1(X,H^1(F))\to\check H^3(X,F)
\]
\LayerTarget{tohoku:criterion:cech-H2-example-reduction:p1}
it is enough \textit{to give a case in which $H^2(X,F)\neq 0$ and
$H^2(X,F)\to\check H^1(X,H^1(F))$ is an isomorphism.}

\TohokuContinuation{178}
